\documentclass[11pt,letter]{article}
\usepackage[colorlinks=true,linkcolor=blue,urlcolor=cyan]{hyperref}
\usepackage[dvipsnames]{xcolor}
\usepackage[margin=0.7in]{geometry}
\usepackage[most]{tcolorbox}
\usepackage{amsmath}
\usepackage{amssymb}
\usepackage{array}
\usepackage{booktabs}
\usepackage{colortbl}
\usepackage{ebgaramond}
%\usepackage{times}
\usepackage{enumitem}
\setlist{nosep}
\usepackage{fancybox}
\usepackage{fancyhdr}
\usepackage{float}
\usepackage{fontawesome5}
\usepackage{geometry}
\usepackage{graphicx}
\usepackage{inconsolata}
\usepackage{indentfirst}
\usepackage{longtable}
\usepackage{makecell}
\usepackage{marginnote}
\usepackage{microtype}
\usepackage{multirow}
\usepackage{soul}
\usepackage{twemojis}
\usepackage{varwidth}
\usepackage{xcolor}
\usepackage{qrcode}

% tikz
\usepackage{tikz}
\usetikzlibrary{positioning, shapes, backgrounds, calc}

% color
\definecolor{intjpurple}{HTML}{88619A}
%\definecolor{intjpurple}{HTML}{4A148C}

% tip box
\definecolor{tipblue}{RGB}{30, 100, 180}
\definecolor{tipbg}{RGB}{246, 249, 253}
\newtcolorbox{tipbox}[1][Tip]{%
  enhanced,
  breakable,
  sidebyside,
  sidebyside align=top,
  lefthand width=4.8em,              % Width allocated to [Icon + Title]
  sidebyside gap=6pt,                % Space between columns
  segmentation style={draw=none},    % Completely removes the dotted separator line
  colback=tipbg,
  colframe=tipblue,
  boxrule=0pt,
  leftrule=3pt,
  arc=0mm,
  boxsep=0pt,
  top=8pt,
  bottom=8pt,
  left=8pt,
  right=8pt,
  before upper={%
    \leavevmode\textcolor{tipblue}{\textbf{\faLightbulb[regular]\enspace #1:}}%
    \tcblower\ignorespaces
  }
}

% macro
\newcommand{\verified}{\colorbox{green!64}{\textcolor{white}{\checkmark Verified}}}

% toc setting
\setcounter{tocdepth}{2}

\title{It's Not You -- How to Identify and Heal from Narcissistic People\\ Dr. Ramani Durvasula\\ \ul{Reading Group Memo}}
\author{}
\date{\today}

\begin{document}
\maketitle
%\tableofcontents

\section*{About this Reading Group Memo}

\begin{tipbox}
\qrcode{https://cdn.penguin.co.uk/dam-assets/books/9781785045028/9781785045028-sample.pdf}
\quad Scan the QR code for a free PDF sample of this book.
\end{tipbox}

\section*{Preface}

I recognize that we can pull ourselves out of the stories of the narcissistic
people who defined us, silenced us, clipped our wings, taught us our dreams
were grandiose, filled us with shame, and for a time, stole our joy.  That we
can have love stories, success stories, and happiness, yet comprehend that
there will still be dark nights of the soul, and that the shadow of self-doubt
remains a fellow traverler.

\section*{Introduction: How Did We Get Here?}

\begin{quote}
Neutrality helps the oppressor, never the victim. Silence encourages the tormentor, never the tormented.
\flushright -- ELIE WIESEL
\end{quote}

\begin{itemize}

\item There is a proverb that says Until the story of the hunt is told by the
lion, the tale of the hunt will always glorify the hunter. The person who holds
the narrative holds the power.

\item The only thing you need to understand about narcissism is that in almost
all cases this personality pattern was there before you came into the
narcissistic person’s life and it will be there after you leave. These
relationships change you, but out of that shift comes growth, a new
perspective, and greater interpersonal discernment. Recognizing and coming out
of these relationships can be a wake-up call to excavate your authentic self,
dust it off, and take it into the world.

\item I hope that this book shines a light on a simple premise that
narcissistic patterns and behaviors really don’t change and you are not to
blame for another person’s invalidating behavior. I want you to land on a
simple yet profound truth:   It’s not you. 

\item To finally be told, clearly and definitively, that you were never going
to be able to change another person’s behavior. All of this was like turning
the house lights on and the gaslights off.

\end{itemize}

\part{The Narcissistic Relationship}

\section{Clarifying Narcissism}

\begin{quote}
The personality susceptible to the dream of limitless freedom is a personality
also prone, should the dream ever sour, to misanthropy and rage. \flushright --
JONATHAN FRANZEN
\end{quote}

{\it Narcissism} is the word of our time, and yet it’s deeply misunderstood. It
would be easier if narcissistic people were simply mirror-gazing, self-involved
poseurs, but they are much more than that. They are an emotionally abusive
romantic partner who belittles you but with whom you sometimes have fun. A
toxic boss who berates you in front of your colleagues but whose work you
greatly admire. A parent who is jealous of your success but who showed up to
all your soccer games when you were a child. A friend who is forever a victim
and drones on endlessly about what’s going on in their life with little
interest in yours but who has been in your life since you were thirteen. Even
these snapshots fail to capture the complexity that is narcissism. You yourself
have likely had one or more relationships with a narcissistic person—and may
not have even known it.

\subsection{What is Narcissism?}

Narcissism is an interpersonally maladaptive personality style that encompasses
a wide spectrum of traits and behavioral patterns that present in different
ways, from mild to severe, vulnerable to malignant. What separates a
narcissistic person from someone who is self-centered or vain or entitled is
the consistency and sheer number of these traits in one person. Just being
superficial doesn’t mean someone is narcissistic. 

It is also looking at the function of these traits—which is to protect the
narcissistic person.  \ul{Narcissism is about a deep insecurity and fragility
offset by maneuvers like domination, manipulation, and gaslighting, which allow
the narcissistic person to stay in control. The variable empathy and lack of
self-awareness mean that they do not stop to consider the harm their behavior
is creating for other people.} The traits aren’t really the issue, but rather
how these traits translate into consistently harmful behaviors.

Because traits, especially those in a rigid and unaware personality style such
as narcissism, really do not change, the behavior is also unlikely to change.
And because there is such a wide spectrum of narcissism, from mild to severe,
we may be having very different experiences of this personality style in our
relationships. \ul{The middle of this spectrum—with enough bad days to take a
toll and enough good days to keep you hooked—is where many people are stuck,
and that “moderate narcissism” is where we will focus.}

Let’s take a closer look at some of these traits.

\begin{description}

\item[A Need for Narcissistic Supply]
Narcissistic individuals constantly demand external validation—via status,
appearance, or flattery—known as narcissistic supply. When denied this supply,
they sour into resentment and irritability, forcing others to either appease
them or endure their wrath.

\item[Egocentricity]
Narcissistic egocentricity combines selfishness with active devaluation, always
putting their needs first while demeaning your preferences (e.g., choosing the
restaurant while insulting your taste in food).

\item[Consistent Inconsistency]
\ul{Narcissism is consistent. However, there is a consistency to it that can make
it feel inconsistent.}
While their moods seem unpredictable, shifts are systematically tied to external
validation: they remain pleasant when praise and control are abundant, but
because supply quickly goes stale, even minor disruptions trigger rapid rage or
sullenness.

\item[Restlessness]
Driven by a restless pursuit of novelty and excitement—often manifesting as
infidelity, overspending, or frenetic activity—narcissistic people quickly
become bored, disenchanted, or contemptuous when life feels routine.

\item[Delusional Grandiosity]
\ovalbox{\it a defining characteristic}
Grandiosity involves exaggerated beliefs of superiority, unique specialness, and
unlimited success; it is "delusional" because they cling to these claims without
real evidence and despite the harm caused to others.

\item[Shifting Masks]
Because their self-appraisal depends entirely on external control, narcissists
swing unpredictably from charming when things go well to enraged or victimized
when things go wrong, leaving you walking on eggshells.

\item[Entitlement]
\ul{Entitlement is a core pattern of narcissism, and one of the most problematic.
Theories of narcissism suggest that entitlement may be the core pillar of this
personality style and that all other dynamics tie back to it.}
Believing rules only apply to ordinary people, narcissists expect VIP treatment
and freedom from accountability, reacting with explosive anger or tantrums when
held to standard boundaries or inconvenienced.

\item[Overcompensating for Insecurity]
Rather than genuine self-esteem, narcissism is driven by deep-seated inadequacy;
their outward arrogance, entitlement, and charm serve merely as defensive armor
shielding an extraordinarily fragile psyche.

\item[Being Thin-skinned]
Narcissists dish out criticism but retaliate against the mildest critique with
disproportionate rage. Attempting to soothe their underlying insecurity is a
trap, as obvious reassurance is often rejected as a patronizing insult that
exposes their weakness.

\item[Inability to Self-regulate]
Because vulnerability feels intolerably shameful, narcissists cannot manage
their emotions and instead discharge internal turmoil by lashing out. When
feedback triggers an ego injury, they impulsively react with rage,
blame-shifting, and hollow apologies to protect their grandiose facade.

\item[Need for Dominance]
Prioritizing power, status, and control over genuine intimacy, narcissists treat
relationships as tools for personal benefit rather than mutual connection.

\item[Lack of Empathy]
Narcissistic empathy is hollow, performative, and transactional, deployed
mainly to gain advantage or manage appearances. Because their support depends on
feeling supplied, they may comfort you on a good day yet dismiss the exact same
struggle as "whining" when their supply is low.

\item[Contempt for Others]
Narcissists need people for validation but deeply resent this dependency,
reacting with hostility or passive-aggressive digs toward others' needs and
vulnerabilities, which mirror their own hidden insecurities.

\item[Projection of Shame]
To protect their grandiose ideal and avoid shame, narcissists unconsciously
project their own unacceptable behaviors onto others—often accusing you of the
exact deceit or harm they are perpetrating themselves.

\item[Being Incredibly Charming]
Acting as skilled social chameleons, narcissists use exceptional charm and
charisma—often mistaken by society for ambition or confidence—to camouflage
toxic traits and draw you close before their harmful behavior surfaces.

\end{description}

\subsection{The Continuum of Narcissism}

\begin{enumerate}

\item The moderate narcissist offers enough good days to keep you invested and
enough bad days that hurt you and leave you utterly confused.

\item Moderate narcissists have just enough insight to know that their behavior is
not okay—but not enough regulation, mindfulness, or empathy to stop themselves.

\item They might compliment you in a meeting in front of colleagues then verbally
eviscerate you behind closed doors in their office. This two-faced
mask-on-mask-off behavior is a trademark of the moderate narcissist.

\end{enumerate}

\subsection{The Different Types of Narcissism}

Although sharing fundamental traits, narcissism manifests across diverse types,
hybrids, and continuums of severity that extend well beyond the commonly
recognized grandiose presentation.

\begin{description}

\item[Grandiose]
Grandiose narcissism masks deep-seated inadequacy behind charismatic arrogance
and fantasy, reacting to setbacks with rage and blame that leave others
exhausted and confused.

\item[Vulnerable]
Vulnerable (or covert) narcissists mask their underlying grandiosity and
entitlement behind chronic resentment, social awkwardness, and perceived
victimhood, externalizing failure by blaming unfair circumstances while
devaluing the success of others.

\item[Communal]
Communal narcissists satisfy their need for validation and a grandiose
self-image by performing visible humanitarian acts, spiritual leadership, or
charitable deeds to win public praise, often while remaining abusive,
indignant, or neglectful behind closed doors.

\item[Self-righteous]
Self-righteous narcissists project a rigid, black-and-white moral superiority
by engineering strictly controlled routines and sneering at anyone who deviates
from their exacting standards of work, discipline, or lifestyle.

\item[Neglectful]
Neglectful narcissists exhibit profound emotional detachment and total
disinterest in close relationships, leaving partners feeling completely
invisible, disregarded, and isolated.

\item[Malignant]
A malignant narcissist represents the dangerous crossroads of narcissism,
psychopathy, sadism, and Machiavellianism—operating as a cruel, paranoid, and
vindictive bully whose aggression and need for domination stem from an
underlying sense of inadequacy and anxiety that pure psychopaths lack.

\end{description}

\subsection{The Battle over Narcissism vs. Narcissistic Personality Disorder}

I understand the fear of labeling: it reduces the complexity of a person to a
single word.

When a person is usually superficially charming but enempathetic, entitled, and
disrespectful, then on a bad day is really nasty, doesn't apologize, and blames
you, that is more likely to be narcissism.

NPD is a paradoxical illness that may harm the people the narcissistic person
interacts with more than it harms the narcissistic individuals themselves.

\subsection{Myths about Narcissism}



\end{document}
